% Discussion 4.1 --- reinforcing findings: the acoustic environment.
\begin{table}[htbp]
\centering
\caption{Space-type findings that reinforce prior intuition: the acoustic environment.}
\label{tab:insights-acoustic}
\footnotesize
\renewcommand{\arraystretch}{1.3}
\begin{tabularx}{\textwidth}{@{}>{\raggedright\arraybackslash\bfseries}p{0.30\textwidth} >{\raggedright\arraybackslash}X@{}}
\toprule
\normalfont\textbf{Key insight} & \textbf{Detail} \\
\midrule
Commute spaces are the noisiest & Transportation and outdoor responses carry the highest \emph{a lot} noise shares ($\approx 25\%$ and $22\%$), against only about 3\% at home. \\
Home and offices are the quietest settings & Both sit near a 62\,dB median sound level, the lowest among the measured spaces. \\
Perceived loudness tracks the setting & Louder spaces attract more \emph{a lot} noise reports, so subjective loudness follows the physical environment. \\
Traffic defines commute noise & Traffic is the dominant reported noise source in transportation ($\approx 57\%$) and outdoors ($\approx 39\%$). \\
Talking defines indoor noise & Talking dominates the acoustic complaints of classrooms ($\approx 81\%$) and offices ($\approx 69\%$). \\
\bottomrule
\end{tabularx}
\end{table}
